% ===========================================================================
\section{Deney protokolü}\label{sec:protokol}
% ===========================================================================

Mimariler arası bir karşılaştırma ancak koşullar eşitse anlamlıdır. Bu bölüm, bütün sonuçların dayandığı kuralları
tanımlar. Kurallar kodda sabittir (\code{sw/zoo/configs.py}, \code{train.py}, \code{experiments.py}).

\paragraph{Parametre bütçesi.} Her mimari, aynı parametre sayısına göre boyutlandırılır. Transformerlerde katman sayısı
veri setine göre sabittir (names'te 4, TinyStories'te 5). Baş boyutu 8'dir; bu, motorun dikkat şeridi sayısıdır. Genişlik
$C$ bütçe kademesinden gelir. MLP'nin gizli genişliği ise parametre sayısı bütçeye en yakın olacak şekilde aranır.
GQA kullanan modellerde $H_{kv}=H/2$'dir. MLP, RNN, GRU ve BoW modellerinde gömme boyutu kademeden gelir, gizli genişlik
aynı kuralla aranır. Bu kural bütün modelleri bütçenin $\pm\%3$'ü içine yerleştirir. Bütçeler names için 12k, 50k ve 200k,
TinyStories için 260k parametredir; sonuncusu Karpathy'nin hazır stories260K modelinin boyutudur. Bu modelin
yapılandırması (5 katman, $C=64$, 8 baş, 4 KV başı, gizli 172) kuralımızın Llama için ürettiği yapılandırmayla birebir
aynıdır.

\paragraph{Başlatma.} Her model, kaynak deposundaki başlatma şemasıyla başlatılır (\code{models.source\_init}).
makemore modelleri PyTorch'un varsayılanını kullanır (makemore da öyle yapar). GPT-2 ve Llama~2, $\mathcal{N}(0, 0{,}02)$
ile ve artık izdüşümleri $0{,}02/\sqrt{2L}$ ile başlatılır; microgpt bütün matrisleri $\mathcal{N}(0, 0{,}08)$ ile
başlatır. nanochat'te gömme $\mathcal{N}(0, 0{,}8)$, çıkış $\mathcal{N}(0, 0{,}001)$ ile başlatılır ve izdüşümler sıfırdır.
Özellik merdiveninin bütün adımları tek bir şema (GPT-2/llama2.c tarzı) kullanır; böylece adımlar arasında yalnızca
mimari değişir.

\paragraph{Eğitim.} AdamW ($\beta=(0{,}9;\,0{,}99)$, ağırlık sönümü 0,01), 500 adım ısınma, tepe değerin 1/30'una kosinüs
azalma, gradyan kırpma 1,0 ve yığın 64 kullanılır. names'te her model 20\,000 adım eğitilir. Öğrenme oranı her mimari
için $\{10^{-3},\;3\cdot10^{-3},\;10^{-2}\}$ kümesinden \emph{doğrulama} kümesinde seçilir. Bunun için names verisi
30\,033 eğitim, 1\,000 doğrulama ve 1\,000 test ismine bölünmüştür; test kümesi önceki sürümlerdekiyle aynıdır. 50k
bütçede seçilen oranla üç tohum (0, 1, 2) eğitilir ve sonuçlar test kümesinde ortalama $\pm$ standart sapma olarak
verilir. Hesap bütçesi nedeniyle 12k bütçede yalnızca oran taraması (tohum 0), 200k bütçede ise 50k'da seçilen oranla
tek tohum çalıştırılır. TinyStories'te her model tek tohum ve tek oranla ($3\cdot10^{-3}$) eğitilir. Bunun için 49
eğitim parçasından rastgele 256 tokenlık pencereler çekilir. Doğrulama ve test pencereleri ise ayrılmış 0 numaralı
parçadan, hikâye başlangıçlarından seçilir.

\paragraph{Özellik merdiveni.} Mimari yeniliklerin her birini ayrı ölçmek için tek bir yapılandırılabilir transformer
(\code{Ladder}) kullanılır. Her adımda tek bir tasarım tercihi değişir:
\begin{enumerate}[label=L\arabic*:, start=0]
  \item GPT-2 bloğu: LayerNorm, öğrenilmiş konum, GELU ve MHA.
  \item LayerNorm yerine RMSNorm.
  \item Öğrenilmiş konum yerine RoPE.
  \item GELU MLP yerine SwiGLU.
  \item MHA yerine GQA.
  \item Bias kaldırılır ve gömme paylaşılır; bu, llama2.c'nin bloğudur.
\end{enumerate}
L6 ve L7 ise nanochat'e uzanan bir yan daldır: L6, L4'te SwiGLU yerine ReLU$^2$ kullanır; L7, L6'ya QK-norm ekler. Her
adımda MLP genişliği bütçe kuralıyla yeniden seçilir. Böylece kalitedeki ve maliyetteki her fark tek bir fikre
bağlanabilir.

\paragraph{Değerlendirme hattı.} Her eğitim koşusu otomatik olarak aynı hattan geçer (\code{pipeline.py}):
\begin{enumerate}
  \item Program derlenir ve INT8'e çevrilir.
  \item names'te test kümesinin tamamında (1\,000 isim), TinyStories'te 64 test penceresinde float ve INT8 kayıp ölçülür.
  \item Model dışa aktarılır ve bir doğrulama motorunda RTL ile çalıştırılır. names'te test kümesinin tamamı, TinyStories'te
  4 tam pencere kullanılır; her logit karşılaştırılır.
\end{enumerate}
Doğrulama motoru, bütün modelleri sığdıracak boyutta tek bir Verilator derlemesidir; modeller \code{ld\_*} portundan
yüklenir. Çevrim sayıları model başına derlemeyle birebir aynıdır.
